New technologies have enabled firms to elicit granular behavioral data from consumers in exchange for lower prices and better experiences. This data can mitigate asymmetric information and moral hazard, but it may also increase firms' market power if kept proprietary. We study a voluntary monitoring program by a major U.S. auto insurer, in which drivers accept short-term tracking in exchange for potential discounts on future premiums. Using a proprietary dataset matched with competitor price menus, we document that safer drivers self-select into monitoring, and those who opt-in become 30\% safer while monitored. Using an equilibrium model of consumer choice and firm pricing for insurance and monitoring, we find that the monitoring program generates large profit and welfare gains. However, large demand frictions hurt monitoring adoption, forcing the firm to offer large discounts to induce opt-in while tempering the profitability of heavy ex-post rent extraction.
A counterfactual data portability policy would further reduce the firm's incentive to elicit monitoring data, leading to less monitoring and lower consumer welfare in equilibrium.